\documentclass[10pt,a4paper]{amsart}
\usepackage{amsmath,amssymb,amsthm,graphicx}
% Retained preamble for the Stacks Project chapter snapshot in this directory.
% The chapter's own preamble.tex is a copy of the shared Stacks Project preamble; the block
% below is sliced MECHANICALLY, line for line, from it. Nothing was added, removed or
% redefined; the wrapper class, the xr-hyper/hyperref wiring and the externaldocument list
% are omitted because the assembler supplies the document class.
%
% lines taken from preamble.tex: 3, 5, 9, 11, 12, 156, 157, 158, ... (51 lines total)
%
% package coverage restored from the chapter's own preamble (2026-10-13): xy, xy-2cell

\usepackage[all]{xy}
\xyoption{2cell}
\UseAllTwocells

\newenvironment{reference}{\comment}{\endcomment}
\newenvironment{slogan}{\comment}{\endcomment}
\newenvironment{history}{\comment}{\endcomment}
\theoremstyle{plain}
\newtheorem{theorem}[subsection]{Theorem}
\newtheorem{proposition}[subsection]{Proposition}
\newtheorem{lemma}[subsection]{Lemma}
\theoremstyle{definition}
\newtheorem{definition}[subsection]{Definition}
\newtheorem{example}[subsection]{Example}
\newtheorem{exercise}[subsection]{Exercise}
\newtheorem{situation}[subsection]{Situation}
\theoremstyle{remark}
\newtheorem{remark}[subsection]{Remark}
\newtheorem{remarks}[subsection]{Remarks}
\numberwithin{equation}{subsection}
\def\lim{\mathop{\mathrm{lim}}\nolimits}
\def\colim{\mathop{\mathrm{colim}}\nolimits}
\def\Spec{\mathop{\mathrm{Spec}}}
\def\Hom{\mathop{\mathrm{Hom}}\nolimits}
\def\Ext{\mathop{\mathrm{Ext}}\nolimits}
\def\SheafHom{\mathop{\mathcal{H}\!\mathit{om}}\nolimits}
\def\SheafExt{\mathop{\mathcal{E}\!\mathit{xt}}\nolimits}
\def\Sch{\mathit{Sch}}
\def\Mor{\mathop{\mathrm{Mor}}\nolimits}
\def\Ob{\mathop{\mathrm{Ob}}\nolimits}
\def\Sh{\mathop{\mathit{Sh}}\nolimits}
\def\NL{\mathop{N\!L}\nolimits}
\def\CH{\mathop{\mathrm{CH}}\nolimits}
\def\proetale{{pro\text{-}\acute{e}tale}}
\def\etale{{\acute{e}tale}}
\def\QCoh{\mathit{QCoh}}
\def\Ker{\mathop{\mathrm{Ker}}}
\def\Im{\mathop{\mathrm{Im}}}
\def\Coker{\mathop{\mathrm{Coker}}}
\def\Coim{\mathop{\mathrm{Coim}}}
\DeclareMathSymbol{\boxtimes}{\mathbin}{AMSa}{"02}
\def\QCohstack{\mathcal{QC}\!\mathit{oh}}
\def\Cohstack{\mathcal{C}\!\mathit{oh}}
\def\Spacesstack{\mathcal{S}\!\mathit{paces}}
\def\Quotfunctor{\mathrm{Quot}}
\def\Hilbfunctor{\mathrm{Hilb}}
\def\Curvesstack{\mathcal{C}\!\mathit{urves}}
\def\Polarizedstack{\mathcal{P}\!\mathit{olarized}}
\def\Complexesstack{\mathcal{C}\!\mathit{omplexes}}
\def\Pic{\mathop{\mathrm{Pic}}\nolimits}
\def\Picardstack{\mathcal{P}\!\mathit{ic}}
\def\Picardfunctor{\mathrm{Pic}}
\def\Deformationcategory{\mathcal{D}\!\mathit{ef}}
\usepackage{cleveref}
\crefname{originallemma}{original lemma}{original lemmas}
\Crefname{originallemma}{Original lemma}{Original lemmas}
\makeatletter
\newlabel{definition-pushforward}{{0FTP}{source}{}{sourcepointer-1}{}}
\newlabel{definition-pushforward@cref}{{[originallemma][0][]0FTP}{[1][0][]source}{}{}{}}
\makeatother
\crefname{originallemma}{original lemma}{original lemmas}
\Crefname{originallemma}{Original lemma}{Original lemmas}
\makeatletter
\newlabel{lemma-compose-pushforward}{{0FTU}{source}{}{sourcepointer-2}{}}
\newlabel{lemma-compose-pushforward@cref}{{[originallemma][0][]0FTU}{[1][0][]source}{}{}{}}
\makeatother
\crefname{originallemma}{original lemma}{original lemmas}
\Crefname{originallemma}{Original lemma}{Original lemmas}
\makeatletter
\newlabel{section-functoriality-dg}{{0FRT}{source}{}{sourcepointer-3}{}}
\newlabel{section-functoriality-dg@cref}{{[originallemma][0][]0FRT}{[1][0][]source}{}{}{}}
\makeatother
\crefname{originallemma}{original lemma}{original lemmas}
\Crefname{originallemma}{Original lemma}{Original lemmas}
\makeatletter
\newlabel{sites-cohomology-section-unbounded}{{07A5}{source}{}{sourcepointer-4}{}}
\newlabel{sites-cohomology-section-unbounded@cref}{{[originallemma][0][]07A5}{[1][0][]source}{}{}{}}
\makeatother
\title{Stacks, Differential Graded Sheaves: lemma-pushforward-agrees}
\author{The Stacks Project Authors}
\date{stacks-project:sdga.tex:local-snapshot-20261012; GNU FDL 1.2 or later}
\begin{document}
\maketitle
\noindent\textit{Editorial note.} Statement span is lines 5377-5396 of sources/stacks-sdga/pinned-original-source.tex; proof span is lines 5398-5448, the complete author proof (49 body lines).\par\medskip
\noindent\textit{Editorial note.} Navigation stubs: 4 cited-but-not-included labels are supplied as non-load-bearing pointers with their pinned tags where the snapshot tag index has them.\par\medskip
\noindent\textit{Editorial note.} extra\_packages is empty on purpose: the retained preamble already loads xy with options and a second declaration clashes in the assembler wrapper.\par\medskip
\noindent\textit{Editorial note.} Text is retained verbatim, including the author's own typos and placeholders.\par\medskip
\noindent\textit{Editorial note.} Queue index 11 of queue\_sdga.json (pb=2091, stmt 5377-5396, proof 5398-5448).\par\medskip
\noindent\textit{Allowed standard prerequisites (editorial).}\begin{itemize}
\item Explicit prerequisite boundary. The cited inputs are prerequisites; the argument in the retained span is not replaced by any cited result.
\end{itemize}\medskip
\ifdefined\hypertarget\hypertarget{sourcepointer-1}{}\fi
\noindent\textit{Source reference.} Original lemma 0FTP: sources/stacks-sdga/pinned-original-source.tex, source lines 5135--5165. Cited by the retained statement or proof but not part of the collected spans; supplied only as a navigation pointer. The argument does not depend on its content.\par\medskip
\ifdefined\hypertarget\hypertarget{sourcepointer-2}{}\fi
\noindent\textit{Source reference.} Original lemma 0FTU: sources/stacks-sdga/pinned-original-source.tex, source lines 5324--5338. Cited by the retained statement or proof but not part of the collected spans; supplied only as a navigation pointer. The argument does not depend on its content.\par\medskip
\ifdefined\hypertarget\hypertarget{sourcepointer-3}{}\fi
\noindent\textit{Source reference.} Original lemma 0FRT: sources/stacks-sdga/pinned-original-source.tex, source lines 2137--2225. Cited by the retained statement or proof but not part of the collected spans; supplied only as a navigation pointer. The argument does not depend on its content.\par\medskip
\ifdefined\hypertarget\hypertarget{sourcepointer-4}{}\fi
\noindent\textit{Source reference.} Original lemma 07A5: sources/stacks-sdga/pinned-original-source.tex, source lines 5448--5448. Cited by the retained statement or proof but not part of the collected spans; supplied only as a navigation pointer. The argument does not depend on its content.\par\medskip
\begin{lemma}
\label{lemma-pushforward-agrees}
Let $(f, f^\sharp) : (\Sh(\mathcal{C}), \mathcal{O}_\mathcal{C})
\to (\Sh(\mathcal{D}), \mathcal{O}_\mathcal{D})$
be a morphism of ringed topoi. Let $\mathcal{A}$ be a differential
graded $\mathcal{O}_\mathcal{C}$-algebra. Let $\mathcal{B}$ be a
differential graded $\mathcal{O}_\mathcal{D}$-algebra. Let
$\varphi : \mathcal{B} \to f_*\mathcal{A}$ be a homomorphism
of differential graded $\mathcal{O}_\mathcal{D}$-algebras.
The diagram
$$
\xymatrix{
D(\mathcal{A}, \text{d}) \ar[d]_{Rf_*} \ar[rr]_{forget} & &
D(\mathcal{O}_\mathcal{C}) \ar[d]^{Rf_*} \\
D(\mathcal{B}, \text{d}) \ar[rr]^{forget} & &
D(\mathcal{O}_\mathcal{D})
}
$$
commutes.
\end{lemma}

\begin{proof}
Besides identifying some categories, this lemma follows immediately
from Lemma \ref{lemma-compose-pushforward}.

\medskip\noindent
We may view $(\mathcal{O}_\mathcal{C}, 0)$ as a differential graded
$\mathcal{O}_\mathcal{C}$-algebra by placing $\mathcal{O}_\mathcal{C}$
in degree $0$ and endowing it with the zero differential. It is clear
that we have
$$
\textit{Mod}(\mathcal{O}_\mathcal{C}, 0) =
\text{Comp}(\mathcal{O}_\mathcal{C})
\quad\text{and}\quad
D(\mathcal{O}_\mathcal{C}, 0) = D(\mathcal{O}_\mathcal{C})
$$
Via this identification the forgetful functor
$\textit{Mod}(\mathcal{A}, \text{d}) \to
\text{Comp}(\mathcal{O}_\mathcal{C})$
is the ``pushforward'' $\text{id}_{\mathcal{C}, *}$
defined in Section \ref{section-functoriality-dg}
corresponding to the identity morphism
$\text{id}_\mathcal{C} : (\mathcal{C}, \mathcal{O}_\mathcal{C}) \to
(\mathcal{C}, \mathcal{O}_\mathcal{C})$ of ringed topoi and the
map $(\mathcal{O}_\mathcal{C}, 0) \to (\mathcal{A}, \text{d})$
of differential graded $\mathcal{O}_\mathcal{C}$-algebras.
Since $\text{id}_{\mathcal{C}, *}$ is exact, we immediately see that
$$
R\text{id}_{\mathcal{C}, *} = forget :
D(\mathcal{A}, \text{d}) \longrightarrow
D(\mathcal{O}_\mathcal{C}, 0) = D(\mathcal{O}_\mathcal{C})
$$
The exact same reasoning shows that
$$
R\text{id}_{\mathcal{D}, *} = forget :
D(\mathcal{B}, \text{d}) \longrightarrow
D(\mathcal{O}_\mathcal{D}, 0) = D(\mathcal{O}_\mathcal{D})
$$
Moreover, the construction of
$Rf_* : D(\mathcal{O}_\mathcal{C}) \to D(\mathcal{O}_\mathcal{D})$
of Cohomology on Sites, Section \ref{sites-cohomology-section-unbounded}
agrees with the construction of
$Rf_* : D(\mathcal{O}_\mathcal{C}, 0) \to D(\mathcal{O}_\mathcal{D}, 0)$
in Definition \ref{definition-pushforward}
as both functors are defined as the right derived extension of
pushforward on underlying complexes of modules.
By Lemma \ref{lemma-compose-pushforward} we see that both
$Rf_* \circ R\text{id}_{\mathcal{C}, *}$ and
$R\text{id}_{\mathcal{D}, *} \circ Rf_*$ are the derived functors
of $f_* \circ forget = forget \circ f_*$ and hence equal
by uniqueness of adjoints.
\end{proof}

\end{document}
